\section{Vector Calculus}
\label{sec:vector-calculus}

\noindent\textbf{Maps and fields.}\quad
\[
f:\R^m\to\R^n,\qquad
\mathbf{x}=x^i\mathbf{e}_i,\qquad
\mathbf{e}_i\cdot\mathbf{e}_j=\delta_{ij},
\qquad i=1,\ldots,n.
\]
Curves: $m=1$. Scalar fields: $\phi:\R^n\to\R$. Vector fields:
$\mathbf{F}:\R^n\to\R^n$. Repeated Cartesian indices are summed.
\src{Tong VC \S 0}

\subsection{Curves}

\noindent\textbf{Parameterized curve and derivative.}
\[
C:\R\to\R^n,\qquad t\mapsto \mathbf{x}(t),\qquad
\mathbf{x}(t+\delta t)-\mathbf{x}(t)
=\dot{\mathbf{x}}(t)\,\delta t+O(\delta t^2).
\]
\[
\dot{\mathbf{x}}=\frac{d\mathbf{x}}{dt}
=\lim_{\delta t\to0}\frac{\delta\mathbf{x}}{\delta t},
\qquad
\mathbf{x}=x^i(t)\mathbf{e}_i\Rightarrow
\dot{\mathbf{x}}=\dot{x}^i(t)\mathbf{e}_i
\quad(\dot{\mathbf{e}}_i=0).
\]
\[
\frac{d}{dt}(f\mathbf{g})=\dot f\,\mathbf{g}+f\dot{\mathbf{g}},
\qquad
\frac{d}{dt}(\mathbf{g}\cdot\mathbf{h})
=\dot{\mathbf{g}}\cdot\mathbf{h}+\mathbf{g}\cdot\dot{\mathbf{h}},
\]
\[
\frac{d}{dt}(\mathbf{g}\times\mathbf{h})
=\dot{\mathbf{g}}\times\mathbf{h}+\mathbf{g}\times\dot{\mathbf{h}}
\quad(\R^3).
\]

\noindent\textbf{Tangent and regularity.}
\[
\dot{\mathbf{x}}(t)\in T_{\mathbf{x}(t)}C,\qquad
\text{regular parameterization: }\dot{\mathbf{x}}(t)\ne0.
\]
The tangent direction is geometric up to orientation; $|\dot{\mathbf{x}}|$ depends on
the parameter. Piecewise smooth $C=C_1+C_2+\cdots$ has tangents away from joins.

\noindent\textbf{Arc length.}
\[
\delta s=|\delta\mathbf{x}|+O(|\delta\mathbf{x}|^2)
=|\dot{\mathbf{x}}\,\delta t|+O(\delta t^2),
\qquad
\frac{ds}{dt}=\pm\left|\frac{d\mathbf{x}}{dt}\right|.
\]
\[
s(t)=\int_{t_0}^{t}|\dot{\mathbf{x}}(t')|\,dt'
\quad(t>t_0).
\]
For $\tau=\tau(t)$ smooth, invertible, $d\tau/dt>0$,
\[
\frac{d\mathbf{x}}{dt}
=\frac{d\mathbf{x}}{d\tau}\frac{d\tau}{dt},
\qquad
\int_{t_0}^{t}|\dot{\mathbf{x}}|\,dt
=\int_{\tau_0}^{\tau}\left|\frac{d\mathbf{x}}{d\tau'}\right|\,d\tau'.
\]
Natural parameter:
\[
\mathbf{t}(s)\equiv\frac{d\mathbf{x}}{ds},\qquad |\mathbf{t}|=1.
\]

\noindent\textbf{Curvature, normal, torsion.}\quad
For $C\subset\R^3$ parameterized by arc length and $\kappa\ne0$,
\[
\kappa(s)\equiv
\left|\frac{d^2\mathbf{x}}{ds^2}\right|
=\left|\frac{d\mathbf{t}}{ds}\right|,
\qquad
\mathbf{n}\equiv
\frac{1}{\kappa}\frac{d^2\mathbf{x}}{ds^2}
=\frac{1}{\kappa}\frac{d\mathbf{t}}{ds}.
\]
\[
\mathbf{t}\cdot\mathbf{t}=1\Rightarrow\mathbf{t}\cdot\mathbf{n}=0,
\qquad
\kappa_{\rm circle}=R^{-1}.
\]
The span of $\{\mathbf{t},\mathbf{n}\}$ is the osculating plane. Define
\[
\mathbf{b}\equiv\mathbf{t}\times\mathbf{n},\qquad
|\mathbf{b}|=1,\qquad
\frac{d\mathbf{b}}{ds}=-\tau(s)\mathbf{n}.
\]
Frenet--Serret:
\[
\boxed{
\frac{d\mathbf{t}}{ds}=\kappa\mathbf{n},\qquad
\frac{d\mathbf{b}}{ds}=-\tau\mathbf{n},\qquad
\frac{d\mathbf{n}}{ds}=\tau\mathbf{b}-\kappa\mathbf{t}}
\]
\[
\frac{d\mathbf{t}}{ds}=\kappa(\mathbf{b}\times\mathbf{t}),
\qquad
\frac{d\mathbf{b}}{ds}=-\tau(\mathbf{b}\times\mathbf{t}).
\]
$\kappa(s),\tau(s)$ plus initial orientation determine $C$ up to rigid motion.
\src{Tong VC \S 1.1}

\noindent\textbf{Scalar line integral.}
\[
\phi:\R^n\to\R,\qquad
\int_C\phi\,ds\equiv\int_{s_a}^{s_b}\phi(\mathbf{x}(s))\,ds.
\]
For $t_a<t_b$,
\[
\boxed{\int_C\phi\,ds
=\int_{t_a}^{t_b}\phi(\mathbf{x}(t))|\dot{\mathbf{x}}(t)|\,dt.}
\]
Orientation independent; $\int_C ds$ is length.

\noindent\textbf{Vector line integral.}
\[
\boxed{\int_C\mathbf{F}\cdot d\mathbf{x}
\equiv\int_{t_a}^{t_b}
\mathbf{F}(\mathbf{x}(t))\cdot\dot{\mathbf{x}}(t)\,dt}
\]
with limits chosen by the orientation of $C$. Hence
\[
\int_{-C}\mathbf{F}\cdot d\mathbf{x}=
-\int_C\mathbf{F}\cdot d\mathbf{x},\qquad
\oint_C\mathbf{F}\cdot d\mathbf{x}=\text{circulation}.
\]
If $C=C_1+C_2+\cdots$,
\[
\int_C\mathbf{F}\cdot d\mathbf{x}
=\int_{C_1}\mathbf{F}\cdot d\mathbf{x}
+\int_{C_2}\mathbf{F}\cdot d\mathbf{x}+\cdots .
\]

\noindent\textbf{Gradient and conservative fields.}
\[
\partial_i\phi\equiv\frac{\partial\phi}{\partial x^i},
\qquad
\nabla\phi\equiv\frac{\partial\phi}{\partial x^i}\mathbf{e}_i,
\qquad
D_{\hat{\mathbf{n}}}\phi=\hat{\mathbf{n}}\cdot\nabla\phi.
\]
$D_{\hat{\mathbf{n}}}\phi$ is maximal for $\hat{\mathbf{n}}\parallel\nabla\phi$.
\[
\mathbf{F}\text{ conservative}\Longleftrightarrow \mathbf{F}=\nabla\phi.
\]
For $C:a\to b$,
\[
\boxed{\mathbf{F}=\nabla\phi\Rightarrow
\int_C\mathbf{F}\cdot d\mathbf{x}=\phi(\mathbf{b})-\phi(\mathbf{a})}
\]
since
\[
\int_C\nabla\phi\cdot d\mathbf{x}
=\int_{t_a}^{t_b}\frac{\partial\phi}{\partial x^i}
\frac{dx^i}{dt}\,dt
=\int_{t_a}^{t_b}\frac{d}{dt}\phi(\mathbf{x}(t))\,dt.
\]
Equivalently, modulo global regularity of the domain and potential,
\[
\boxed{\mathbf{F}\text{ conservative}
\Longleftrightarrow
\oint_C\mathbf{F}\cdot d\mathbf{x}=0\text{ for all closed }C.}
\]
Local test:
\[
\mathbf{F}=\nabla\phi\Rightarrow
F_i=\partial_i\phi,\qquad
\partial_jF_i=\partial_iF_j.
\]
The symmetry $\partial_iF_j=\partial_jF_i$ is necessary and locally sufficient for
regular fields on simple domains.

\noindent\textbf{Exact differential.}
\[
d\phi=\frac{\partial\phi}{\partial x^i}\,dx^i
=\nabla\phi\cdot d\mathbf{x},
\qquad
\mathbf{F}\cdot d\mathbf{x}=d\phi
\Longleftrightarrow \mathbf{F}=\nabla\phi.
\]

\noindent\textbf{Work and potential energy.}
\[
m\ddot{\mathbf{x}}=\mathbf{F}(\mathbf{x}),\qquad
K=\frac12m|\dot{\mathbf{x}}|^2,
\qquad
K(t_2)-K(t_1)=\int_C\mathbf{F}\cdot d\mathbf{x}.
\]
If $\mathbf{F}=-\nabla V$,
\[
K(t_2)-K(t_1)=-V(\mathbf{x}(t_2))+V(\mathbf{x}(t_1)),
\qquad
\boxed{E=K+V=\text{constant}.}
\]

\noindent\textbf{Global subtlety.}\quad
On $\R^2\setminus\{0\}$,
\[
\mathbf{F}(x,y)=
\left(-\frac{y}{x^2+y^2},\,\frac{x}{x^2+y^2}\right),
\qquad
\frac{\partial F_x}{\partial y}=\frac{\partial F_y}{\partial x}
=\frac{y^2-x^2}{(x^2+y^2)^2}.
\]
Locally $\mathbf{F}=\nabla\tan^{-1}(y/x)$, but the potential is not globally
continuous and $\mathbf{F}$ is singular at $0$. For
$\mathbf{x}(t)=(R\cos t,R\sin t)$,
\[
\oint_C\mathbf{F}\cdot d\mathbf{x}
=\int_0^{2\pi}
\left[-\frac{\sin t}{R}(-R\sin t)+
\frac{\cos t}{R}(R\cos t)\right]dt
=2\pi.
\]
Closed-loop vanishing requires field and potential data to be well behaved on the
relevant domain, including the region enclosed by the loop.
\src{Tong VC \S 1.2-\S 1.3}

\subsection{Surfaces and Volumes}

\noindent\textbf{Area integrals.}\quad For $D\subset\R^2$,
\[
\int_D \phi\,dA,\qquad dA=dx\,dy,\qquad
\operatorname{Area}(D)=\int_D dA .
\]
Successive integrations encode the shape in limits:
\[
\int_D\phi(x,y)\,dA
=\int_a^bdy\int_{x_1(y)}^{x_2(y)}dx\,\phi(x,y)
=\int_c^ddx\int_{y_1(x)}^{y_2(x)}dy\,\phi(x,y).
\]
Fubini: order interchangeable for sufficiently regular $\phi,D$.

\noindent\textbf{Coordinate changes in \texorpdfstring{$\R^2$}{R2}.}
For a smooth invertible map $(u,v)\mapsto(x(u,v),y(u,v))$,
\[
\boxed{\int_D dx\,dy\,\phi(x,y)
=\int_{D'}du\,dv\,|J(u,v)|\,\phi(x(u,v),y(u,v))}
\]
\[
J=\frac{\partial(x,y)}{\partial(u,v)}
=\det\begin{pmatrix}
\partial_u x & \partial_v x\\
\partial_u y & \partial_v y
\end{pmatrix},
\qquad
\delta A=|J|\,\delta u\,\delta v.
\]
Plane polar:
\[
x=\rho\cos\varphi,\qquad y=\rho\sin\varphi,\qquad
\rho\ge0,\quad\varphi\in[0,2\pi),
\]
\[
J=\rho,\qquad \boxed{dA=\rho\,d\rho\,d\varphi},\qquad
\int_0^\infty e^{-x^2/2}\,dx=\sqrt{\frac{\pi}{2}}.
\]

\noindent\textbf{Volume integrals and coordinate changes.}
\[
\int_V\phi(\mathbf{x})\,dV
=\lim_{\delta V\to0}\sum_n\phi(\mathbf{x}_n)\delta V,
\qquad
\operatorname{Vol}(V)=\int_V dV.
\]
For $(u,v,w)\mapsto(x,y,z)$,
\[
\boxed{dV=dx\,dy\,dz=|J|\,du\,dv\,dw},
\qquad
J=\frac{\partial(x,y,z)}{\partial(u,v,w)}
=\det\begin{pmatrix}
\partial_u x & \partial_v x & \partial_w x\\
\partial_u y & \partial_v y & \partial_w y\\
\partial_u z & \partial_v z & \partial_w z
\end{pmatrix}.
\]
Spherical polar:
\[
x=r\sin\theta\cos\varphi,\quad
y=r\sin\theta\sin\varphi,\quad
z=r\cos\theta,
\]
\[
r\in[0,\infty),\quad\theta\in[0,\pi],\quad\varphi\in[0,2\pi),
\qquad
\boxed{dV=r^2\sin\theta\,dr\,d\theta\,d\varphi}.
\]
\[
\int_{B_R} f(r)\,dV=4\pi\int_0^R r^2f(r)\,dr,
\qquad
\operatorname{Vol}(B_R)=\frac{4\pi R^3}{3}.
\]
Cylindrical polar:
\[
x=\rho\cos\varphi,\qquad y=\rho\sin\varphi,\qquad z=z,
\]
\[
\rho\in[0,\infty),\quad \varphi\in[0,2\pi),\quad z\in\R,
\qquad
\boxed{dV=\rho\,d\rho\,d\varphi\,dz}.
\]
Sphere with cylinder removed, $x^2+y^2+z^2\le R^2$, $x^2+y^2\ge s^2$:
\[
\operatorname{Vol}
=\int_0^{2\pi}d\varphi\int_s^R d\rho\,\rho
\int_{-\sqrt{R^2-\rho^2}}^{\sqrt{R^2-\rho^2}}dz
=\boxed{\frac{4\pi}{3}(R^2-s^2)^{3/2}}.
\]

\noindent\textbf{Vector-valued volume integrals and center of mass.}
\[
M=\int_V\rho(\mathbf{x})\,dV,\qquad
\mathbf{X}=\frac{1}{M}\int_V\mathbf{x}\,\rho(\mathbf{x})\,dV.
\]
Uniform solid upper hemisphere:
\[
M=\frac{2\pi}{3}\rho R^3,\qquad \mathbf{X}=(0,0,3R/8).
\]
In $\R^n$, for $x^i=x^i(u^1,\ldots,u^n)$,
\[
\int_M f(x^i)\,dx^1\cdots dx^n
=\int_{M'}f(x^i(u))\,|J|\,du^1\cdots du^n,
\qquad
J=\det\left(\frac{\partial x^i}{\partial u^a}\right).
\]

\noindent\textbf{Surfaces in \texorpdfstring{$\R^3$}{R3}.}
\[
S=\{\mathbf{x}\in\R^3:F(x,y,z)=0\},
\qquad
\mathbf{x}=\mathbf{x}(u,v):D\subset\R^2\to\R^3.
\]
Level-set normal:
\[
\mathbf{n}=\pm\frac{\nabla F}{|\nabla F|},\qquad
\text{regular when }\nabla F\ne0\text{ on }S.
\]
Parametric tangents and normal:
\[
\mathbf{x}_u=\frac{\partial\mathbf{x}}{\partial u},\quad
\mathbf{x}_v=\frac{\partial\mathbf{x}}{\partial v},\quad
\mathbf{n}\sim\mathbf{x}_u\times\mathbf{x}_v,\quad
\mathbf{x}_u\times\mathbf{x}_v\ne0.
\]
Examples:
\[
F=x^2+y^2+z^2-R^2,\quad \nabla F=2(x,y,z),
\]
\[
F=x^2+y^2-z^2-R^2,\quad \nabla F=2(x,y,-z).
\]

\noindent\textbf{Boundaries, orientation, closedness.}
\[
\partial S=\text{piecewise smooth closed boundary curve},\qquad
\partial V=\text{boundary surface of }V,\qquad
\boxed{\partial^2=0}.
\]
An orientable surface admits a globally smooth unit normal. Reversing orientation:
\[
\mathbf{n}\mapsto-\mathbf{n}.
\]
A bounded surface lies in a finite ball; a closed surface is bounded and has no boundary.

\noindent\textbf{Scalar surface integral.}
\[
dS=|\mathbf{x}_u\times\mathbf{x}_v|\,du\,dv,
\qquad
\boxed{\int_S\phi(\mathbf{x})\,dS
=\int_Ddu\,dv\,|\mathbf{x}_u\times\mathbf{x}_v|\,
\phi(\mathbf{x}(u,v)).}
\]
Orientation independent; $\operatorname{Area}(S)=\int_SdS$. Under
$(u,v)\mapsto(\tilde u,\tilde v)$,
\[
\mathbf{x}_u\times\mathbf{x}_v
=\frac{\partial(\tilde u,\tilde v)}{\partial(u,v)}
\mathbf{x}_{\tilde u}\times\mathbf{x}_{\tilde v}.
\]
Sphere of radius $R$:
\[
\mathbf{x}(\theta,\varphi)
=R(\sin\theta\cos\varphi,\sin\theta\sin\varphi,\cos\theta)=R\mathbf{e}_r,
\]
\[
\mathbf{x}_\theta=R\mathbf{e}_\theta,\qquad
\mathbf{x}_\varphi=R\sin\theta\,\mathbf{e}_\varphi,\qquad
\mathbf{x}_\theta\times\mathbf{x}_\varphi
=R^2\sin\theta\,\mathbf{e}_r.
\]
\[
\boxed{dS=R^2\sin\theta\,d\theta\,d\varphi},\qquad
A_{\rm cap}(0\le\theta\le\alpha)=\boxed{2\pi R^2(1-\cos\alpha)}.
\]

\noindent\textbf{Flux.}
\[
d\mathbf{S}=\mathbf{n}\,dS
=(\mathbf{x}_u\times\mathbf{x}_v)\,du\,dv.
\]
\[
\boxed{\int_S\mathbf{F}\cdot d\mathbf{S}
=\int_Ddu\,dv\,
(\mathbf{x}_u\times\mathbf{x}_v)\cdot\mathbf{F}(\mathbf{x}(u,v))}
\]
Changing orientation flips the sign. Fluid rate through $S$:
\[
\text{rate}=\int_S\mathbf{F}\cdot d\mathbf{S}.
\]
For $\mathbf{F}=(-x,0,z)$ on the radius-$R$ cap $0\le\theta\le\alpha$,
\[
\int_S\mathbf{F}\cdot d\mathbf{S}
=\boxed{\pi R^3\cos\alpha\,\sin^2\alpha}.
\]

\noindent\textbf{Gaussian curvature and Gauss--Bonnet data.}
\[
\kappa_{\min}\le\kappa\le\kappa_{\max},\qquad
\boxed{K=\kappa_{\min}\kappa_{\max}}.
\]
Gauss theorem egregium: $K$ is intrinsic. Geodesic triangle:
\[
\boxed{\theta_1+\theta_2+\theta_3
=\pi+\int_DK\,dS}.
\]
Closed orientable surface of genus $g$:
\[
\boxed{\int_SK\,dS=4\pi(1-g)},\qquad
K_{S_R^2}=R^{-2},\qquad \int_{S^2_R}K\,dS=4\pi.
\]
\src{Tong VC \S 2}

\subsection{Grad, Div, Curl}

\noindent\textbf{Gradient.}
\[
\nabla\phi=\frac{\partial\phi}{\partial x^i}\mathbf{e}_i,\qquad
\phi(\mathbf{x}+\mathbf{h})
=\phi(\mathbf{x})+\mathbf{h}\cdot\nabla\phi+O(|\mathbf{h}|^2).
\]
Along $\mathbf{x}(t)$,
\[
\frac{d}{dt}\phi(\mathbf{x}(t))
=\frac{\partial\phi}{\partial x^i}\frac{dx^i}{dt}
=\nabla\phi\cdot\dot{\mathbf{x}}.
\]
Radial identities, $r^2=x^ix^i$, $\mathbf{r}=x^i\mathbf{e}_i=r\hat{\mathbf{r}}$:
\[
\frac{\partial r}{\partial x^i}=\frac{x^i}{r},\qquad
\nabla r^p=pr^{p-1}\hat{\mathbf{r}},\qquad
\nabla\cdot\mathbf{r}=n,\qquad
\nabla\times\mathbf{r}=0\quad(n=3).
\]
\[
\phi=-\frac1r\Rightarrow \nabla\phi=\frac{\hat{\mathbf{r}}}{r^2}.
\]

\noindent\textbf{Operators.}
\[
\nabla=\mathbf{e}_i\frac{\partial}{\partial x^i},\qquad
\nabla\cdot\mathbf{F}=\frac{\partial F_i}{\partial x^i}.
\]
In $\R^3$,
\[
\nabla\times\mathbf{F}
=\epsilon_{ijk}\frac{\partial F_j}{\partial x^i}\mathbf{e}_k
=\begin{pmatrix}
\partial_2F_3-\partial_3F_2\\
\partial_3F_1-\partial_1F_3\\
\partial_1F_2-\partial_2F_1
\end{pmatrix}.
\]
\[
\epsilon_{ijk}\epsilon_{ilm}
=\delta_{jl}\delta_{km}-\delta_{jm}\delta_{kl}.
\]
Examples:
\[
\mathbf{F}=(x^2,0,0):\quad\nabla\cdot\mathbf{F}=2x,\quad\nabla\times\mathbf{F}=0,
\]
\[
\mathbf{F}=(y,-x,0):\quad\nabla\cdot\mathbf{F}=0,\quad\nabla\times\mathbf{F}=(0,0,-2),
\]
\[
\mathbf{F}=\frac{\hat{\mathbf{r}}}{r^2}:\quad
\nabla\cdot\mathbf{F}=0,\quad\nabla\times\mathbf{F}=0\quad(r\ne0),
\qquad
\nabla\cdot\frac{\hat{\mathbf{r}}}{r^2}=4\pi\delta^3(\mathbf{x}).
\]

\noindent\textbf{Linearity and product identities.}
\[
\nabla(\alpha\phi+\psi)=\alpha\nabla\phi+\nabla\psi,
\quad
\nabla\cdot(\alpha\mathbf{F}+\mathbf{G})
=\alpha\nabla\cdot\mathbf{F}+\nabla\cdot\mathbf{G},
\]
\[
\nabla\times(\alpha\mathbf{F}+\mathbf{G})
=\alpha\nabla\times\mathbf{F}+\nabla\times\mathbf{G}.
\]
\[
\nabla(\phi\psi)=\phi\nabla\psi+\psi\nabla\phi,
\]
\[
\nabla\cdot(\phi\mathbf{F})
=(\nabla\phi)\cdot\mathbf{F}+\phi\nabla\cdot\mathbf{F},
\qquad
\nabla\times(\phi\mathbf{F})
=(\nabla\phi)\times\mathbf{F}+\phi\nabla\times\mathbf{F}.
\]
\[
\nabla\cdot(\mathbf{F}\times\mathbf{G})
=(\nabla\times\mathbf{F})\cdot\mathbf{G}
-\mathbf{F}\cdot(\nabla\times\mathbf{G}).
\]
With $(\mathbf{F}\cdot\nabla)\equiv F_i\partial_i$,
\[
\nabla(\mathbf{F}\cdot\mathbf{G})
=\mathbf{F}\times(\nabla\times\mathbf{G})
+\mathbf{G}\times(\nabla\times\mathbf{F})
+(\mathbf{F}\cdot\nabla)\mathbf{G}
+(\mathbf{G}\cdot\nabla)\mathbf{F},
\]
\[
\nabla\times(\mathbf{F}\times\mathbf{G})
=(\nabla\cdot\mathbf{G})\mathbf{F}
-(\nabla\cdot\mathbf{F})\mathbf{G}
+(\mathbf{G}\cdot\nabla)\mathbf{F}
-(\mathbf{F}\cdot\nabla)\mathbf{G}.
\]

\noindent\textbf{Irrotational, solenoidal, Helmholtz.}
On fields defined everywhere on $\R^3$,
\[
\boxed{\nabla\times\mathbf{F}=0\Longleftrightarrow\mathbf{F}=\nabla\phi}
\]
and equivalently
\[
\nabla\times\mathbf{F}=0
\Longleftrightarrow
\mathbf{F}=\nabla\phi
\Longleftrightarrow
\oint_C\mathbf{F}\cdot d\mathbf{x}=0.
\]
\[
\mathbf{F}=\nabla\phi\Rightarrow
\nabla\times\mathbf{F}
=\epsilon_{ijk}\partial_i\partial_j\phi\,\mathbf{e}_k=0.
\]
Solenoidal:
\[
\nabla\cdot\mathbf{F}=0,\qquad
\boxed{\nabla\cdot\mathbf{F}=0\Longleftrightarrow\mathbf{F}=\nabla\times\mathbf{A}}
\quad(\R^3,\text{ global regularity}).
\]
Gauge freedom:
\[
\mathbf{A}\sim\mathbf{A}+\nabla\chi.
\]
One explicit potential for solenoidal $\mathbf{F}$:
\[
\mathbf{A}=
\left(
\int_{z_0}^{z}F_y(x,y,z')\,dz',
\int_{x_0}^{x}F_z(x',y,z_0)\,dx'
-\int_{z_0}^{z}F_x(x,y,z')\,dz',
0\right).
\]
Helmholtz form:
\[
\mathbf{F}=\nabla\phi+\nabla\times\mathbf{A}.
\]

\noindent\textbf{Laplacian.}
\[
\nabla^2\equiv\nabla\cdot\nabla=\partial_i\partial_i,\qquad
\nabla^2_{\R^3}=
\frac{\partial^2}{\partial x^2}
+\frac{\partial^2}{\partial y^2}
+\frac{\partial^2}{\partial z^2}.
\]
\[
\nabla\times(\nabla\times\mathbf{F})
=\nabla(\nabla\cdot\mathbf{F})-\nabla^2\mathbf{F}.
\]
Physics templates:
\[
\nabla\cdot\mathbf{E}=\frac{\rho}{\epsilon_0},\quad
\nabla\times\mathbf{E}=-\frac{\partial\mathbf{B}}{\partial t},\quad
\nabla\cdot\mathbf{B}=0,\quad
\nabla\times\mathbf{B}=\mu_0\left(\mathbf{J}
+\epsilon_0\frac{\partial\mathbf{E}}{\partial t}\right).
\]
\[
\frac{\partial T}{\partial t}=D\nabla^2T,\qquad
i\hbar\frac{\partial\psi}{\partial t}
=-\frac{\hbar^2}{2m}\nabla^2\psi+V\psi.
\]
\src{Tong VC \S 3.1-\S 3.2}

\noindent\textbf{Orthogonal curvilinear coordinates.}
\[
d\mathbf{x}
=\frac{\partial\mathbf{x}}{\partial u}du
+\frac{\partial\mathbf{x}}{\partial v}dv
+\frac{\partial\mathbf{x}}{\partial w}dw,
\qquad
\frac{\partial\mathbf{x}}{\partial u}\cdot
\left(\frac{\partial\mathbf{x}}{\partial v}\times
\frac{\partial\mathbf{x}}{\partial w}\right)\ne0.
\]
\[
\frac{\partial\mathbf{x}}{\partial u}=h_u\mathbf{e}_u,\quad
\frac{\partial\mathbf{x}}{\partial v}=h_v\mathbf{e}_v,\quad
\frac{\partial\mathbf{x}}{\partial w}=h_w\mathbf{e}_w,\quad
\mathbf{e}_u\times\mathbf{e}_v=\mathbf{e}_w.
\]
\[
d\mathbf{x}=h_u\mathbf{e}_u\,du+h_v\mathbf{e}_v\,dv+h_w\mathbf{e}_w\,dw,
\qquad
d\mathbf{x}^2=h_u^2du^2+h_v^2dv^2+h_w^2dw^2.
\]
Cylindrical:
\[
\mathbf{x}=(\rho\cos\phi,\rho\sin\phi,z),\quad
\mathbf{e}_\rho=(\cos\phi,\sin\phi,0),\quad
\mathbf{e}_\phi=(-\sin\phi,\cos\phi,0),\quad
\mathbf{e}_z=\hat{\mathbf{z}},
\]
\[
h_\rho=1,\qquad h_\phi=\rho,\qquad h_z=1.
\]
Spherical:
\[
\mathbf{x}=(r\sin\theta\cos\phi,r\sin\theta\sin\phi,r\cos\theta),
\]
\[
\mathbf{e}_r=(\sin\theta\cos\phi,\sin\theta\sin\phi,\cos\theta),
\quad
\mathbf{e}_\theta=(\cos\theta\cos\phi,\cos\theta\sin\phi,-\sin\theta),
\quad
\mathbf{e}_\phi=(-\sin\phi,\cos\phi,0),
\]
\[
h_r=1,\qquad h_\theta=r,\qquad h_\phi=r\sin\theta.
\]

\noindent\textbf{Grad, div, curl, Laplacian.}
\[
\boxed{\nabla f
=\frac1{h_u}\frac{\partial f}{\partial u}\mathbf{e}_u
+\frac1{h_v}\frac{\partial f}{\partial v}\mathbf{e}_v
+\frac1{h_w}\frac{\partial f}{\partial w}\mathbf{e}_w}
\]
\[
\boxed{\nabla\cdot\mathbf{F}
=\frac1{h_uh_vh_w}
\left[\partial_u(h_vh_wF_u)
+\partial_v(h_uh_wF_v)
+\partial_w(h_uh_vF_w)\right]}
\]
\[
\boxed{\nabla\times\mathbf{F}
=\frac1{h_uh_vh_w}
\begin{vmatrix}
h_u\mathbf{e}_u & h_v\mathbf{e}_v & h_w\mathbf{e}_w\\
\partial_u & \partial_v & \partial_w\\
h_uF_u & h_vF_v & h_wF_w
\end{vmatrix}}
\]
where derivatives act on the third row only.
\[
\boxed{\nabla^2f
=\frac1{h_uh_vh_w}
\left[
\partial_u\left(\frac{h_vh_w}{h_u}\partial_u f\right)
+\partial_v\left(\frac{h_uh_w}{h_v}\partial_v f\right)
+\partial_w\left(\frac{h_uh_v}{h_w}\partial_w f\right)
\right].}
\]
Cylindrical:
\[
\nabla f=f_\rho\hat{\bm{\rho}}+\frac1{\rho}f_\phi\hat{\bm{\phi}}+f_z\hat{\mathbf{z}},
\]
\[
\nabla\cdot\mathbf{F}
=\frac1{\rho}\partial_\rho(\rho F_\rho)
+\frac1{\rho}\partial_\phi F_\phi+\partial_zF_z,
\]
\[
\nabla\times\mathbf{F}
=\left(\frac1{\rho}\partial_\phi F_z-\partial_zF_\phi\right)\hat{\bm{\rho}}
+\left(\partial_zF_\rho-\partial_\rho F_z\right)\hat{\bm{\phi}}
+\frac1{\rho}\left[\partial_\rho(\rho F_\phi)-\partial_\phi F_\rho\right]\hat{\mathbf{z}},
\]
\[
\nabla^2f=
\frac1{\rho}\partial_\rho(\rho\partial_\rho f)
+\frac1{\rho^2}\partial_\phi^2f+\partial_z^2f.
\]
Spherical:
\[
\nabla f=f_r\hat{\mathbf{r}}+\frac1r f_\theta\hat{\bm{\theta}}
+\frac1{r\sin\theta}f_\phi\hat{\bm{\phi}},
\]
\[
\nabla\cdot\mathbf{F}
=\frac1{r^2}\partial_r(r^2F_r)
+\frac1{r\sin\theta}\partial_\theta(\sin\theta F_\theta)
+\frac1{r\sin\theta}\partial_\phi F_\phi,
\]
\[
\nabla\times\mathbf{F}
=\frac1{r\sin\theta}
\left[\partial_\theta(\sin\theta F_\phi)-\partial_\phi F_\theta\right]\hat{\mathbf{r}}
\]
\[
\hspace{2em}
+\frac1r\left[\frac1{\sin\theta}\partial_\phi F_r-\partial_r(rF_\phi)\right]\hat{\bm{\theta}}
+\frac1r\left[\partial_r(rF_\theta)-\partial_\theta F_r\right]\hat{\bm{\phi}},
\]
\[
\nabla^2f=
\frac1{r^2}\partial_r(r^2\partial_r f)
+\frac1{r^2\sin\theta}\partial_\theta(\sin\theta\,\partial_\theta f)
+\frac1{r^2\sin^2\theta}\partial_\phi^2f.
\]
\src{Tong VC \S 3.3}

\subsection{Integral Theorems}

\noindent\textbf{Divergence theorem.}\quad
For bounded $V\subset\R^3$ with piecewise smooth closed boundary
$S=\partial V$ and outward $d\mathbf{S}=\mathbf{n}\,dS$,
\[
\boxed{\int_V\nabla\cdot\mathbf{F}\,dV
=\int_{\partial V}\mathbf{F}\cdot d\mathbf{S}.}
\]
Local flux definition:
\[
\boxed{\nabla\cdot\mathbf{F}(\mathbf{x})
=\lim_{V\to0}\frac{1}{V}
\int_{\partial V}\mathbf{F}\cdot d\mathbf{S}.}
\]
Small box:
\[
\int_{\partial(\delta V)}\mathbf{F}\cdot d\mathbf{S}
=\left(\frac{\partial F_x}{\partial x}
+\frac{\partial F_y}{\partial y}
+\frac{\partial F_z}{\partial z}\right)\delta x\,\delta y\,\delta z+o(\delta V).
\]
Tiling cancels interior face fluxes pairwise; only $\partial V$ remains.

\noindent\textbf{Two-dimensional version and scalar corollary.}
For $D\subset\R^2$, $C=\partial D$, outward normal $\mathbf{n}$,
\[
\boxed{\int_D\nabla\cdot\mathbf{F}\,dA
=\int_C\mathbf{F}\cdot\mathbf{n}\,ds.}
\]
For scalar $\phi$,
\[
\boxed{\int_V\nabla\phi\,dV=\int_{\partial V}\phi\,d\mathbf{S}}
\]
from Gauss applied to $\mathbf{F}=\phi\mathbf{a}$ with constant $\mathbf{a}$.
\src{Tong VC \S 4.1}

\noindent\textbf{Continuity equation.}
For density $\rho(\mathbf{x},t)$ and current $\mathbf{J}(\mathbf{x},t)$,
\[
Q_V(t)=\int_V\rho\,dV,\qquad
\boxed{\frac{\partial\rho}{\partial t}+\nabla\cdot\mathbf{J}=0},
\]
\[
\boxed{\frac{dQ_V}{dt}
=-\int_{\partial V}\mathbf{J}\cdot d\mathbf{S}}.
\]
If $\mathbf{J}\to0$ sufficiently fast at infinity, then
\[
Q_{\rm tot}=\int_{\R^3}\rho\,dV,\qquad
\frac{dQ_{\rm tot}}{dt}=0.
\]
Electric current:
\[
I=\int_S\mathbf{J}\cdot d\mathbf{S}.
\]
Fluid mass conservation, $\mathbf{J}=\rho\mathbf{u}$:
\[
\frac{\partial\rho}{\partial t}+\nabla\cdot(\rho\mathbf{u})=0,
\qquad
\rho=\text{const}\Rightarrow \boxed{\nabla\cdot\mathbf{u}=0}.
\]

\noindent\textbf{Diffusion from conservation.}
For energy density $E=cT$ and heat current $\mathbf{J}=-\kappa\nabla T$,
\[
\frac{\partial E}{\partial t}+\nabla\cdot\mathbf{J}=0
\Rightarrow
\boxed{\frac{\partial T}{\partial t}=D\nabla^2T},
\qquad
D=\frac{\kappa}{c}.
\]

\noindent\textbf{Predator-prey / Bendixson-Dulac template.}
Lotka--Volterra:
\[
\dot w=w(-\alpha+c),\qquad
\dot c=c(\beta-w),\qquad \alpha,\beta>0,
\]
\[
(w,c)_*=(\beta,\alpha),\qquad
\frac{dw}{dc}=\frac{w(c-\alpha)}{c(\beta-w)},
\qquad
\boxed{\beta\log w-w+\alpha\log c-c=\text{const}.}
\]
With competition:
\[
\dot w=w(-\alpha+c-\mu w),\qquad
\dot c=c(\beta-w-\nu c),\qquad \mu,\nu>0,
\]
\[
(1+\mu\nu)w_*=\beta-\nu\alpha,\qquad
(1+\mu\nu)c_*=\alpha+\mu\beta.
\]
Let $\mathbf{a}=(w,c)$, $\dot{\mathbf{a}}=\mathbf{p}$, $\nabla=(\partial_w,\partial_c)$.
A closed orbit would imply, for any Dulac multiplier $b$,
\[
0=\oint b\,\mathbf{p}\cdot\mathbf{n}\,dt
=\int_D\nabla\cdot(b\mathbf{p})\,dA.
\]
Choose $b=(wc)^{-1}$:
\[
\boxed{\nabla\cdot\left(\frac{\mathbf{p}}{wc}\right)
=-\frac{\mu}{c}-\frac{\nu}{w}<0\quad(w,c>0).}
\]
Therefore no closed phase-plane orbits if $\mu\ne0$ or $\nu\ne0$.
\src{Tong VC \S 4.2}

\noindent\textbf{Green theorem in the plane.}
For $C=\partial A$ positively oriented,
\[
\boxed{\int_A\left(\frac{\partial Q}{\partial x}
-\frac{\partial P}{\partial y}\right)dA
=\oint_C P\,dx+Q\,dy.}
\]
Equivalence with 2d Gauss: set $\mathbf{F}=(Q,-P)$; for
$C:\mathbf{x}(s)=(x(s),y(s))$,
\[
\mathbf{n}=(y',-x'),\qquad
\mathbf{F}\cdot\mathbf{n}\,ds=Q\,dy+P\,dx.
\]
For holes: outer boundary counterclockwise, inner boundaries clockwise.
\src{Tong VC \S 4.3}

\noindent\textbf{Stokes theorem.}
For oriented $S\subset\R^3$ with boundary $C=\partial S$,
\[
\boxed{\int_S(\nabla\times\mathbf{F})\cdot d\mathbf{S}
=\oint_C\mathbf{F}\cdot d\mathbf{x}.}
\]
Orientation: if $\mathbf{n}$ points toward the observer, $C$ is counterclockwise;
equivalently $\mathbf{t}\times\mathbf{n}$ points outward from $S$.
\[
\boxed{\mathbf{n}\cdot(\nabla\times\mathbf{F})
=\lim_{A\to0}\frac{1}{A}\oint_C\mathbf{F}\cdot d\mathbf{x}.}
\]
Rigid rotation:
\[
\mathbf{u}=\bm{\omega}\times\mathbf{x}
\Rightarrow \nabla\times\mathbf{u}=2\bm{\omega}.
\]
If $\nabla\times\mathbf{F}=0$ on a simply connected region, then
$\oint_C\mathbf{F}\cdot d\mathbf{x}=0$ for all closed $C$, so
$\mathbf{F}=\nabla\phi$.

\noindent\textbf{Stokes from Green.}
For $\mathbf{x}=\mathbf{x}(u,v)$ and $C=\partial S$ corresponding to $\partial A$,
\[
\oint_C\mathbf{F}\cdot d\mathbf{x}
=\oint_{\partial A}F_u\,du+F_v\,dv,
\qquad
F_u=\mathbf{F}\cdot\mathbf{x}_u,\quad
F_v=\mathbf{F}\cdot\mathbf{x}_v,
\]
\[
\frac{\partial F_v}{\partial u}
-\frac{\partial F_u}{\partial v}
=(\nabla\times\mathbf{F})\cdot(\mathbf{x}_u\times\mathbf{x}_v).
\]

\noindent\textbf{Magnetic field of an infinite wire.}
\[
\nabla\times\mathbf{B}=\mu_0\mathbf{J},\qquad
\oint_{\partial S}\mathbf{B}\cdot d\mathbf{x}
=\mu_0\int_S\mathbf{J}\cdot d\mathbf{S}
=\mu_0 I.
\]
For a straight wire along $z$,
\[
\boxed{\mathbf{B}(\rho,\phi)=
\frac{\mu_0I}{2\pi\rho}(-\sin\phi,\cos\phi,0)
=\frac{\mu_0I}{2\pi\rho}\hat{\bm{\phi}}},\qquad
\nabla\cdot\mathbf{B}=0.
\]

\noindent\textbf{Coordinate changes revisited.}
For orthogonal coordinates $(u,v,w)$ with scale factors $h_u,h_v,h_w$,
\[
\boxed{\nabla\cdot\mathbf{F}
=\frac{1}{h_uh_vh_w}\left[
\partial_u(h_vh_wF_u)
+\partial_v(h_uh_wF_v)
+\partial_w(h_uh_vF_w)\right]}
\]
\[
\boxed{
\nabla\times\mathbf{F}
=\frac{1}{h_uh_vh_w}
\begin{vmatrix}
h_u\mathbf{e}_u & h_v\mathbf{e}_v & h_w\mathbf{e}_w\\
\partial_u & \partial_v & \partial_w\\
h_uF_u & h_vF_v & h_wF_w
\end{vmatrix}.}
\]
\src{Tong VC \S 4.4}

\subsection{Poisson and Laplace Equations}

\noindent\textbf{Gravity/electrostatics analogy and Gauss law.}
\[
\mathbf{F}_{\rm test}=m\mathbf{g}+q\mathbf{E},\qquad
\mathbf{g}(\mathbf{x})=-\frac{GM}{r^2}\hat{\mathbf{r}},\qquad
\mathbf{E}(\mathbf{x})=\frac{Q}{4\pi\epsilon_0r^2}\hat{\mathbf{r}}.
\]
For a point source enclosed by closed $S$,
\[
\boxed{\int_S\mathbf{g}\cdot d\mathbf{S}=-4\pi GM},
\qquad
\boxed{\int_S\mathbf{E}\cdot d\mathbf{S}=\frac{Q}{\epsilon_0}}.
\]
For densities $\rho_m,\rho_e$,
\[
\boxed{\nabla\cdot\mathbf{g}=-4\pi G\rho_m},
\qquad
\boxed{\nabla\cdot\mathbf{E}=\frac{\rho_e}{\epsilon_0}}.
\]
Spherical symmetry:
\[
\mathbf{g}(r)=-\frac{GM(r)}{r^2}\hat{\mathbf{r}},\qquad
\mathbf{E}(r)=\frac{Q(r)}{4\pi\epsilon_0r^2}\hat{\mathbf{r}}.
\]
Infinite line charge of linear density $\lambda$:
\[
\boxed{\mathbf{E}(r)=\frac{\lambda}{2\pi\epsilon_0r}\hat{\mathbf{r}}.}
\]
In $\R^n$, point-source radial fields scale as $r^{-(n-1)}$.

\noindent\textbf{Potentials and Poisson equations.}
\[
\nabla\times\mathbf{g}=0,\qquad \nabla\times\mathbf{E}=0,\qquad
\mathbf{g}=-\nabla\Phi,\qquad \mathbf{E}=-\nabla\phi.
\]
\[
\boxed{\nabla^2\Phi=4\pi G\rho_m},\qquad
\boxed{\nabla^2\phi=-\frac{\rho_e}{\epsilon_0}}.
\]
Abstract normalization:
\[
\boxed{\nabla^2\psi(\mathbf{x})=-\rho(\mathbf{x})},\qquad
\rho=0\Rightarrow \boxed{\nabla^2\psi=0}.
\]
Linearity: $\nabla^2\psi_1=\nabla^2\psi_2=0\Rightarrow
\nabla^2(\psi_1+\psi_2)=0$.
\src{Tong VC \S 5.1-\S 5.2}

\noindent\textbf{Isotropic harmonic functions.}
In $\R^3$, $\psi=\psi(r)$:
\[
\nabla^2\psi=\frac{1}{r^2}\frac{d}{dr}
\left(r^2\frac{d\psi}{dr}\right)=0
\Rightarrow
\boxed{\psi(r)=\frac{A}{r}+B}\qquad(r\ne0).
\]
In cylindrical symmetry / $\R^2$:
\[
\nabla^2\psi=\frac{1}{r}\frac{d}{dr}
\left(r\frac{d\psi}{dr}\right)=0
\Rightarrow
\boxed{\psi(r)=A\log r+B}\qquad(r\ne0).
\]
In $\R^n$,
\[
\psi(r)\sim
\begin{cases}
r^{2-n}, & n\ne2,\\
\log r, & n=2.
\end{cases}
\]
Derivatives of harmonic functions are harmonic away from singularities:
\[
\psi_{\rm dipole}(\mathbf{x})
=\mathbf{d}\cdot\nabla\frac{1}{r}
=-\frac{\mathbf{d}\cdot\mathbf{x}}{r^3}.
\]

\noindent\textbf{Spherical constant-density source.}
For
\[
\nabla^2\psi=
\begin{cases}
-\rho_0, & r\le R,\\
0, & r>R,
\end{cases}
\qquad
\psi(0)\text{ finite},\quad \psi(r)\to0,\quad
\psi,\psi'\text{ continuous at }R,
\]
\[
\boxed{
\psi(r)=
\begin{cases}
\dfrac{\rho_0}{6}(3R^2-r^2), & r\le R,\\[4pt]
\dfrac{\rho_0R^3}{3r}, & r>R.
\end{cases}}
\]
For gravity in the above normalization, $\Phi=-4\pi G\psi$.

\noindent\textbf{Boundary data and uniqueness.}
On bounded $V$,
\[
\text{Dirichlet: }\psi|_{\partial V}=f,\qquad
\text{Neumann: }\frac{\partial\psi}{\partial n}
=\mathbf{n}\cdot\nabla\psi=g.
\]
Dirichlet data determine a Poisson solution uniquely; pure Neumann data determine it
up to an additive constant. If $\psi=\psi_1-\psi_2$,
\[
\nabla^2\psi=0,\qquad
\int_V\nabla\cdot(\psi\nabla\psi)\,dV
=\int_V|\nabla\psi|^2\,dV
=\int_{\partial V}\psi\,\mathbf{n}\cdot\nabla\psi\,dS.
\]
Homogeneous Dirichlet or Neumann boundary term vanishes, so $\nabla\psi=0$ in $V$.
Neumann compatibility for $\nabla^2\psi=\rho$:
\[
\boxed{\int_V\rho\,dV=\int_{\partial V}g\,dS.}
\]

\noindent\textbf{Green identities.}
\[
\boxed{
\int_V\phi\nabla^2\psi\,dV
=-\int_V\nabla\phi\cdot\nabla\psi\,dV
+\int_{\partial V}\phi\nabla\psi\cdot d\mathbf{S}}
\]
\[
\boxed{
\int_V(\phi\nabla^2\psi-\psi\nabla^2\phi)\,dV
=\int_{\partial V}(\phi\nabla\psi-\psi\nabla\phi)\cdot d\mathbf{S}.}
\]

\noindent\textbf{Harmonic functions.}
\[
\boxed{\nabla^2\psi=0.}
\]
Mean value property: if $B_R(\mathbf{a})\subset V$,
\[
\boxed{\psi(\mathbf{a})
=\frac{1}{4\pi R^2}
\int_{|\mathbf{x}-\mathbf{a}|=R}\psi(\mathbf{x})\,dS.}
\]
Reason:
\[
\frac{d}{dR}\left[
\frac{1}{4\pi R^2}\int_{S_R}\psi\,dS\right]
=\frac{1}{4\pi R^2}\int_{S_R}\nabla\psi\cdot d\mathbf{S}
=\frac{1}{4\pi R^2}\int_{B_R}\nabla^2\psi\,dV=0.
\]
Maximum/minimum principle:
\[
\boxed{\text{nonconstant harmonic }\psi\text{ has no interior maximum or minimum}.}
\]

\noindent\textbf{Green function in \texorpdfstring{$\R^3$}{R3}.}
\[
\delta^3(\mathbf{x})=0\quad(\mathbf{x}\ne0),\qquad
\int_V f(\mathbf{x})\delta^3(\mathbf{x})\,dV=f(\mathbf{0})\quad(0\in V).
\]
\[
\boxed{\nabla^2\left(\frac{1}{|\mathbf{x}|}\right)
=-4\pi\delta^3(\mathbf{x})}.
\]
For localized source decaying sufficiently fast,
\[
\boxed{\psi(\mathbf{x})
=\frac{1}{4\pi}\int
\frac{\rho(\mathbf{x}')}{|\mathbf{x}-\mathbf{x}'|}\,dV'}
\]
solves $\nabla^2\psi=-\rho$.
\src{Tong VC \S 5.2}

\subsection{Tensors}

\noindent\textbf{Transformation law.}
For orthonormal bases related by $e'_i=R_{ij}e_j$,
\[
RR^T=I,\qquad R\in O(n),\qquad \det R=\pm1.
\]
Vector and rank-$p$ tensor components:
\[
\mathbf{x}=x_i e_i=x'_i e'_i,\qquad x'_i=R_{ij}x_j,
\]
\[
\boxed{T'_{i_1\cdots i_p}
=R_{i_1j_1}\cdots R_{i_pj_p}T_{j_1\cdots j_p}.}
\]
Rank $0$: scalar; rank $1$: vector; rank $2$: matrix only if it obeys this law.

\noindent\textbf{Tensors as maps.}
\[
T(a,b,\ldots,c)=T_{i_1\cdots i_p}a_{i_1}b_{i_2}\cdots c_{i_p}
\]
is multilinear and basis-independent. Rank $2$ maps vectors to vectors:
\[
u_i=T_{ij}v_j,\qquad
T'_{ij}=R_{ik}R_{jl}T_{kl},\qquad T'=RTR^T.
\]

\noindent\textbf{Tensor operations and quotient rule.}
\[
(S+T)_{i_1\cdots i_p}=S_{i_1\cdots i_p}+T_{i_1\cdots i_p},
\]
\[
(S\otimes T)_{i_1\cdots i_pj_1\cdots j_q}
=S_{i_1\cdots i_p}T_{j_1\cdots j_q},
\]
\[
\delta_{ij}T_{ijk_1\cdots k_{p-2}},\qquad
\operatorname{tr}T=T_{ii},\qquad
a\cdot b=a_ib_i.
\]
Quotient rule:
\[
v_{i_1\cdots i_p}
=T_{i_1\cdots i_pj_1\cdots j_q}u_{j_1\cdots j_q}
\text{ tensor for all tensor }u
\Rightarrow T\text{ is a tensor}.
\]

\noindent\textbf{Symmetry and invariant tensors.}
\[
T_{\cdots ij\cdots}=\pm T_{\cdots ji\cdots}
\Rightarrow \text{symmetric/antisymmetric property is basis-independent}.
\]
In $\R^3$,
\[
T_{ij}=P_{ij}+\epsilon_{ijk}B_k+\frac13\delta_{ij}Q,
\qquad
P_{ii}=0,\quad Q=T_{ii},\quad
B_k=\frac12\epsilon_{ijk}T_{ij}.
\]
\[
\delta'_{ij}=R_{ik}R_{jl}\delta_{kl}=\delta_{ij},
\]
\[
\epsilon'_{i_1\cdots i_n}
=R_{i_1j_1}\cdots R_{i_nj_n}\epsilon_{j_1\cdots j_n}
=(\det R)\epsilon_{i_1\cdots i_n}.
\]
Thus $\epsilon$ is invariant under $SO(n)$ and changes sign under reflections.
In $\R^3$:
\[
\text{rank }1:0,\qquad
\text{rank }2:\alpha\delta_{ij},\qquad
\text{rank }3:\beta\epsilon_{ijk},
\]
\[
T_{ijkl}^{\rm iso}
=\alpha\delta_{ij}\delta_{kl}
+\beta\delta_{ik}\delta_{jl}
+\gamma\delta_{il}\delta_{jk},
\qquad
\epsilon_{ijk}\epsilon_{ilp}
=\delta_{jl}\delta_{kp}-\delta_{jp}\delta_{kl}.
\]
\src{Tong VC \S 6.1}

\noindent\textbf{Invariant integrals.}
For rotationally invariant $V$ and radial $f(r)$,
\[
T_{i\cdots k}=\int_V f(r)x_i\cdots x_k\,dV
\]
is isotropic. Hence in $\R^3$,
\[
\int_V\rho(r)x_i\,dV=0,
\]
\[
\int_V\rho(r)x_ix_j\,dV
=\frac13\delta_{ij}\int_V\rho(r)r^2\,dV
=\frac{4\pi}{3}\delta_{ij}\int_0^R\rho(r)r^4\,dr.
\]

\noindent\textbf{Tensor fields and vector-gradient decomposition.}
\[
X_{i_1\cdots i_qj_1\cdots j_p}
=\partial_{i_1}\cdots\partial_{i_q}
T_{j_1\cdots j_p}(\mathbf{x}),\qquad
\partial_i\equiv\frac{\partial}{\partial x_i}.
\]
For vector field $F_i$,
\[
\partial_jF_i=P_{ij}+\epsilon_{ijk}B_k+\frac13\delta_{ij}Q,
\]
\[
Q=\partial_iF_i=\nabla\cdot\mathbf{F},\qquad
B_k=\frac12\epsilon_{ijk}\partial_jF_i
=-\frac12(\nabla\times\mathbf{F})_k,
\]
\[
P_{ij}=\frac12(\partial_jF_i+\partial_iF_j)
-\frac13\delta_{ij}\partial_kF_k.
\]

\noindent\textbf{Physical tensors.}
\[
P_i=\alpha_{ij}E_j,\qquad J_i=\sigma_{ij}E_j.
\]
Isotropic $3$d media:
\[
\sigma_{ij}=\sigma\delta_{ij},\qquad
\alpha_{ij}=\alpha\delta_{ij}.
\]
Isotropic $2$d conductor:
\[
\sigma_{ij}=\sigma_{xx}\delta_{ij}+\sigma_{xy}\epsilon_{ij}
=\begin{pmatrix}
\sigma_{xx} & \sigma_{xy}\\
-\sigma_{xy} & \sigma_{xx}
\end{pmatrix}.
\]
\[
I_{ij}=\sum_a m_a\left(|\mathbf{x}_a|^2\delta_{ij}
-(x_a)_i(x_a)_j\right),\qquad \mathbf{L}=I\bm{\omega}.
\]
Continuous body:
\[
\boxed{I_{ij}=\int_V\rho(\mathbf{x})
\left(|\mathbf{x}|^2\delta_{ij}-x_ix_j\right)\,dV.}
\]
Sphere, radius $R$, radial density:
\[
I_{ij}=\frac{8\pi}{3}\delta_{ij}\int_0^R\rho(r)r^4\,dr,
\qquad
\rho=\rho_0:\quad I_{ij}=\frac25MR^2\delta_{ij}.
\]
Cylinder, radius $a$, height $2L$:
\[
I=\operatorname{diag}\left(
M\left(\frac{a^2}{4}+\frac{L^2}{3}\right),
M\left(\frac{a^2}{4}+\frac{L^2}{3}\right),
\frac12Ma^2\right).
\]
Since $I_{ij}=I_{ji}$, principal axes diagonalize $I$:
\[
I'=RIR^T=\operatorname{diag}(I_1,I_2,I_3).
\]
Elasticity and geometry templates:
\[
\sigma_{ij}=C_{ijkl}e_{kl},\qquad R_{ijkl}\text{ (Riemann curvature)}.
\]
\src{Tong VC \S 6.2}

\noindent\textbf{Antisymmetric tensors and integration.}
Generalized cross product:
\[
(a\times b)_{i_1\cdots i_{n-2}}
=\epsilon_{i_1\cdots i_{n-2}jk}a_jb_k.
\]
Hodge-dual-type contraction:
\[
(*T)_{i_1\cdots i_{n-p}}
=\frac1{p!}\epsilon_{i_1\cdots i_{n-p}j_1\cdots j_p}
T_{j_1\cdots j_p}.
\]
A totally antisymmetric $p$-tensor integrates over a $p$-dimensional subspace $M$:
\[
\int_M T_{i_1\cdots i_p}\,dS_{i_1\cdots i_p}.
\]
For $\mathbf{x}=\mathbf{x}(u,v)$,
\[
dS_{ij}
=\frac12\left(
\frac{\partial x_i}{\partial u}\frac{\partial x_j}{\partial v}
-\frac{\partial x_i}{\partial v}\frac{\partial x_j}{\partial u}
\right)du\,dv.
\]

\noindent\textbf{Exterior-derivative-like operator.}
\[
(D\phi)_i=\partial_i\phi,\qquad
(D\mathbf{F})_{ij}=\frac12(\partial_jF_i-\partial_iF_j),
\]
\[
(\nabla\times\mathbf{F})_i
=\epsilon_{ijk}(D\mathbf{F})_{jk}\quad(\R^3).
\]
For antisymmetric $p$-tensor $T$,
\[
(DT)_{i_1\cdots i_{p+1}}
=\frac1{p+1}\sum_{r=1}^{p+1}(-1)^{r+1}
\partial_{i_r}T_{i_1\cdots\widehat{i_r}\cdots i_{p+1}},
\qquad
D^2T=0.
\]

\noindent\textbf{Generalized Stokes theorem.}
\[
\boxed{\int_M(DT)_{i_1\cdots i_{p+1}}\,dS_{i_1\cdots i_{p+1}}
=\int_{\partial M}T_{i_1\cdots i_p}\,dS_{i_1\cdots i_p}.}
\]
Here $\dim M=p+1$ and $\dim\partial M=p$. Special cases: fundamental theorem of
calculus, divergence theorem, ordinary Stokes theorem.
\src{Tong VC \S 6.3}
